\documentclass[12pt,x11names,a4paper]{article}
%Preamble

\usepackage{comment}
\usepackage[utf8]{inputenc}
\usepackage{amsmath}
\usepackage{amssymb}
\usepackage{amsthm}
\usepackage{enumitem}
\usepackage{mathrsfs}
\usepackage{mathtools}
\usepackage{esvect}
\usepackage{xfp}
\usepackage{xparse}
\usepackage{array}
\usepackage{changepage}
\usepackage{pdfpages}
\usepackage[parfill]{parskip}
\usepackage{anyfontsize}
\usepackage{alphalph}
\usepackage{ulem}
\usepackage{cancel}
\usepackage[framemethod=tikz]{mdframed}
\usepackage{xargs}
\usepackage{siunitx}              


\usepackage{algorithm}
\usepackage{algorithmic}
\usepackage{matlab-prettifier}
\usepackage{textcomp}
\usepackage{listings}
\usepackage[dvipsnames,svgnames]{xcolor}
\usepackage{colortbl}
\usepackage{eso-pic}
\usepackage[margin=1.35in]{geometry}
\usepackage{float}
\usepackage[danish]{babel}
\usepackage{tikz}
\usetikzlibrary{matrix,calc}
\usetikzlibrary{arrows}
\usetikzlibrary{arrows.meta}
\usetikzlibrary{decorations.pathreplacing}
\usetikzlibrary{decorations.pathreplacing,calligraphy}
\usetikzlibrary{patterns}
\usetikzlibrary{angles, quotes}
\usepackage{csquotes}
\usepackage{pdfpages}
\usepackage{fancyhdr}
\usepackage{lastpage}
\usepackage{pgfplots}
\usepackage{pgffor}

%Ændre følgende, hvis du vil have en anden bredde på pgfplots
\pgfplotsset{width=10cm,compat = 1.9}




\usepackage{caption}
\usepackage{subcaption}
\makeatletter
% \def\@seccntformat#1{%
%   \expandafter\ifx\csname c@#1\endcsname\c@section\else
%   \csname the#1\endcsname\quad
%   \fi}
 \makeatother
\counterwithin*{equation}{section}
\numberwithin{equation}{section}



%\newtheorem{setn}{Sætning}[section]
%\newtheorem{prop}[defn]{Proposition}
%\newtheorem{lem}[defn]{Lemma}
%\newtheorem{cor}[defn]{Korollar}

\newcommand\Tr{\textnormal{Tr}}
\newcommand\N{\textnormal{N}}
\newcommand{\suco}[2]{\left.#1\right|_{\mathbb{F}_{#2}}}
\renewcommand\qedsymbol{$\blacksquare$}
\theoremstyle{definition}
%\newtheorem{defn}[setn]{Definition}

\newcommand\Ker{\textnormal{Ker}\hspace{0.5mm}}
\renewcommand\Im{\textnormal{Im}\hspace{0.5mm}}
\newcommand\TS{\hspace{0.8mm}}

\usepackage{tcolorbox}
\tcbuselibrary{theorems}

\newmdtheoremenv[
	skipabove = 5pt,
	skipbelow = 5pt,
	roundcorner = 5,
	backgroundcolor = SeaGreen!30,
	linecolor = white, 
	]
	{defn}{ Definition }[section]

\newmdtheoremenv[
	skipabove = 5pt,
	skipbelow = 5pt,
	roundcorner = 5,
	backgroundcolor = RubineRed!20,
	linecolor = white, 
	]
	{setn}[defn]{ Sætning }
	

\newtheorem{eksempel}[defn]{Eksempel}

\newenvironment{exa}
	{\begin{eksempel}}
	{\unskip\hfill $\square$ \end{eksempel}}
\newtheorem{regel}[defn]{Regel}


\usepackage{graphicx}
\usepackage[backend=biber]{biblatex}
\addbibresource{biblio.bib}

\renewcommand{\mod}[1]{\ (\textnormal{mod }{#1})}
\usepackage{hyperref, bookmark}
\usepackage{booktabs}
\usepackage{calc}
\usepackage[nottoc,numbib]{tocbibind} 
\usepackage{multirow,bigdelim}

\renewcommand{\arraystretch}{1.3}
\newcommand\intd{\textnormal{d}}
\newcommand\e{\textnormal{e}}




\newenvironment{opgavetekst}[1]
{
	\begin{minipage}[t]{0.15\textwidth}
	 	\textbf{ #1}
	\end{minipage}
	\begin{minipage}[t]{0.85\textwidth}
		\rule{\textwidth}{0.3mm} \\
}
{ 
\\	\end{minipage}
}

\newenvironment{meretekst}
{
	\begin{minipage}[t]{0.15\textwidth}
	 	\phantom{h}
	\end{minipage}
	\begin{minipage}[t]{0.85\textwidth}
}
{ 
\\	\end{minipage}
}

\newenvironment{delopgave}[1]
{
	\begin{minipage}[t]{0.15\textwidth}
		\phantom{\textbf{Opgave}}		
		\textbf{\symbol{\numexpr96+ #1})}
	\end{minipage}
	\begin{minipage}[t]{0.85\textwidth}

	\begin{minipage}[t]{\textwidth}
}{
\end{minipage}
\phantom{h}\\
\end{minipage}
}


%\newcommand{opgavelinje}[1]{
%\begin{minipage}[t]{0.15\textwidth}
%\textbf{#1}
%\end{minipage} 
%\begin{minipage}[t]{\textwidth}
%\rule{0.81\textwidth}{0.3mm}
%\end{minipage}
%}





%Farver!
\definecolor{NorregGroen}{RGB}{0,105,78}





\newgeometry{margin=2cm}

\pagestyle{fancy}
\fancyhf{}

\rhead{Nørre Gymnasium\\2.i
}
\cfoot{Side \thepage \hspace{1pt} af \pageref{LastPage}}

%Husk at rette modul og dato!
\lhead{Assignment 4\\ Analysis and Approaches SL
}
\chead{December 19 2025
}

\begin{document}

%\includepdf[pages=-]{Forsider/aarsprove_1v.pdf}
\savegeometry{art}

\begin{titlepage}
\newgeometry{margin=0pt}

\begin{minipage}{0.27\textwidth}

\begin{tikzpicture}[overlay]
\fill[top color = NorregGroen!40, bottom color = NorregGroen] (6,10) rectangle (-10,-30);
\end{tikzpicture}
\end{minipage}
\begin{minipage}{0.73\textwidth}
\begin{center}
\phantom{h} \vspace{1cm}\\
\hspace{4cm}
\includegraphics[scale = 1]{Billeder/Norreg.png} \\
\phantom{h} \vspace{5cm}\\
\rule{0.7\textwidth}{0.3mm}\\
\phantom{h}\\
{\fontsize{50}{60}\selectfont Assignment 4}\\
\phantom{h}\\
\rule{0.7\textwidth}{0.3mm}\\
\Large 2025\\
\Large 2.i maSL AA

\end{center}
\end{minipage}
\end{titlepage}
\loadgeometry{art}

%Udfyld afsnit herunder og lav til egen Latex-fil

%Kopier følgende til overskrift:

%\begin{center}
%\Huge
%Aflevering 1
%\end{center}
%\section*{Opgave 1}
%\stepcounter{section}
\begin{center}
%Opgavesætter er delt i to dele:\\
%Delprøve 1 kun med den centralt udmeldte formelsamling.\\
%Delprøve 2 med alle hjælpemidler.
\end{center}

%\section*{Krav til formidling af din besvarelse}

%Ved bedømmelse af helhedsindtrykket af besvarelsen af de enkelte opgaver lægges særlig vægt på følgende fire punkter:
%\begin{itemize}
%\item[$\cdot$] \textbf{Redegørelse og dokumentation for metode} \\
%Besvarelsen skal indeholde en redegørelse for den anvendte løsningsstragegi med dokumentation i form af et passende antal mellemregninger \textit{eller} matematiske forklaringer på metoden, når et matematisk værktøjsprogram anvendes.
%\item[$\cdot$] \textbf{Figurer, grafer og andre illustrationer} \\
%Besvarelsen skal indeholde hensigtsmæssig brug af figurer, grafer og andre illustrationer, og der skal være tydelige henvisninger til brug af disse i den forklarende tekst.
%\item[$\cdot$] \textbf{Notation og layout}\\
%Besvarelsen skal i overensstemmelse med god matematisk skik opstilles med hensigtsmæssig brug af symbolsprog, og med en redegørelse for den matematiske notation, der indføres og anvendes, og som ikke kan henføres stil standardviden.
%\item[$\cdot$] \textbf{Formidling og forklaring}\\
%Besvarelsen af rene matematikopgaver skal indeholde en angivelse af givne oplysninger og korte forklaringer knyttet til den anvendte løsningsstrategi beskrevet med brug af almindelig matematisk notation. 

%Besvarelsen af opgaver, der omhandler matematiske modeller, skal indeholde en kort præsentation af modellens kontekst, herunder betydning af modellens parametre. De enkelte delspørgsmål skal afsluttes med en præcis konklusion præsenteret i et klart sprog i relation til konteksten.
%\end{itemize}

%\newpage

\begin{opgavetekst}{Problem 1}
	The function $f$ is given by
	\begin{align*}
		f(x) = 5x^2 -10x + 9.
	\end{align*}
\end{opgavetekst}
\begin{delopgave}{1}
	Determine an antiderivative of $f$ and show that your function is an antiderivative.
\end{delopgave}

\begin{opgavetekst}{Problem 2}
	The function $g$ is given by
	\begin{align*}
		g(x) = \frac{7}{x} + \frac{3}{5}, \ x> 0.
	\end{align*}
\end{opgavetekst}
\begin{delopgave}{1}
	Determine $\int g(x) dx$.
\end{delopgave}

\begin{opgavetekst}{Problem 3}
	The indeterminate integral of the function $\ln(x)$ is given by
	\begin{align*}
		\int \ln(x) dx = x\ln(x) - x + C
	\end{align*}
\end{opgavetekst}
\begin{delopgave}{1}
	Give a proof of this rule. 
\end{delopgave}


\begin{opgavetekst}{Problem 4}
	The function $h$ is given by
	\begin{align*}
		h(t) = \frac{t^3}{\sqrt{t^4+1}}
	\end{align*}
\end{opgavetekst}
\begin{delopgave}{1}
	Find $\int h(t) dt$.
\end{delopgave}

\begin{opgavetekst}{Problem 5}	
\end{opgavetekst}
\begin{delopgave}{1}
	Calculate the following
	\begin{align*}
		&1) \int_1^3 3x^2+2x-1 dx & &2) \int_{9}^{16} \frac{1}{2\sqrt{x}} dx
	\end{align*}
\end{delopgave}

\begin{opgavetekst}{Problem 6}
	A function $\varphi$ is given by
	\begin{align*}
		\varphi(x) = \frac{e^{4\sin(x)}\cos(x)}{2}
	\end{align*}
\end{opgavetekst}
\begin{delopgave}{1}
	Determine 
	\begin{align*}
		\int_0^\pi \varphi(x) dx
	\end{align*}
\end{delopgave}

\begin{opgavetekst}{Problem 7}
	The functions $f$ and $g$ given by 
	\begin{align*}
		f(x) = -x^2 + 2x +3
	\end{align*}
	and 
	\begin{align*}
		g(x) = 3x-9
	\end{align*}
	bound a region $M$.
\end{opgavetekst}
\begin{delopgave}{1}
	Determine the intersections between $f$ and $g$
\end{delopgave}
\begin{delopgave}{2}
	Determine the roots of $f$
\end{delopgave}
\begin{delopgave}{3}
	Determine the area of $M$. 
\end{delopgave}

\begin{opgavetekst}{Problem 8}
	The evil exercise a) from the book $>:)$. 
\end{opgavetekst}
\begin{delopgave}{1}
	Determine the integral
		\begin{align*}
			\int \frac{x^3}{\sqrt{x^2+1}}dx.
		\end{align*}
\end{delopgave}
\end{document}



